\documentclass[12pt]{article}

\usepackage[T2A]{fontenc}
\usepackage[english, russian]{babel}
\usepackage[letterpaper,top=2cm,bottom=2cm,left=2cm,right=2cm,marginparwidth=1.75cm]{geometry}
\usepackage{amsmath,amsthm,amssymb}
\usepackage{bbm}
\usepackage{hyperref}
\hypersetup{
    colorlinks=true,     
    % linkcolor=blue,
    % pdftitle={Overleaf Example},
    % pdfpagemode=FullScreen,
    }
% \usepackage[normalem]{ulem}
% \usepackage{graphicx}
% \usepackage{subcaption}
% \usepackage{pgfplots}
% \pgfplotsset{width=10cm,compat=1.18}
% \usepgfplotslibrary{external}
% \tikzexternalize[prefix=tikz/]
% \tikzexternaldisable
% \usepackage{xcolor}
\usepackage{mathabx}
\usepackage{tikz}
\usetikzlibrary {arrows.meta}
\usetikzlibrary {decorations.pathmorphing}
\usetikzlibrary {decorations.markings}
\usetikzlibrary{bending}

\tikzset{
    pics/cross/.style ={
        code = {
            \draw (-2pt,-2pt) -- (2pt,2pt) (-2pt,2pt) -- (2pt,-2pt);
        }
    },
    pics/atp/.style={
        code = {
            \draw (-2pt,-2pt) -- (2pt,2pt) (-2pt,2pt) -- (2pt,-2pt);
            \draw[-{Stealth[length=7pt, bend]}] (0,0.2) arc[start angle=90, end angle=270, radius=0.2cm];
            \draw[-{Stealth[length=7pt, bend]}] (0,-0.2) arc[start angle=-90, end angle=90, radius=0.2cm];
        }
    }
}

\newcommand{\N}{\mathbb{N}}
\newcommand{\Z}{\mathbb{Z}}
\newcommand{\R}{\mathbb{R}}
\newcommand{\Q}{\mathbb{Q}}
\renewcommand{\C}{\mathbb{C}}
\newcommand{\F}{\mathbb{F}}
\newcommand{\ind}{\mathbbm{1}}
\renewcommand{\P}{\mathbb{P}}
\newcommand{\E}{\mathbb{E}}
\newcommand{\D}{\mathbb{D}}
\newcommand{\vol}{\operatorname{vol}}
\newcommand{\Int}{\operatorname{Int}}
\newcommand{\Cl}{\operatorname{Cl}}
\newcommand{\Fr}{\operatorname{Fr}}
\newcommand{\Gal}{\operatorname{Gal}}
\newcommand{\id}{\operatorname{id}}
\newcommand{\Aut}{\operatorname{Aut}}
\newcommand{\End}{\operatorname{End}}
\newcommand{\Res}{\operatorname{Res}}
\renewcommand{\Re}{\operatorname{Re}}
\renewcommand{\Im}{\operatorname{Im}}
\renewcommand{\ind}{\operatorname{ind}}
\newcommand{\inter}{\cap}
\newcommand{\union}{\cup}
\newcommand{\Ra}{\Rightarrow}
\newcommand{\La}{\Leftarrow}
\renewcommand{\iff}{\Leftrightarrow}
\newcommand{\todo}[1]{\colorbox{red}{#1}}

\newcounter{cnt}

\newtheoremstyle{defnote}
{\topsep}
{\topsep}
{}
{}
{\bfseries}
{.}
{ }
{\thmname{#1}\thmnumber{ #2}\thmnote{ (#3)}}

\newtheoremstyle{thmnote}
{\topsep}
{\topsep}
{\itshape}
{}
{\bfseries}
{.}
{ }
{\thmname{#1}\thmnumber{ #2}\thmnote{ (#3)}}

\theoremstyle{thmnote}
\newtheorem{theorem}[cnt]{Теорема}
\newtheorem{lemma}[cnt]{Лемма}
\newtheorem{proposition}[cnt]{Предложение}

\theoremstyle{remark}
\newtheorem*{remark}{Замечание}

\theoremstyle{defnote}
\newtheorem{definition}[cnt]{Определение}



\begin{document}

\title{Матанализ, 4 семестр}
\author{}
\date{}
\maketitle

\section*{Комплексный анализ}

Будем рассматривать функции из $\C$ в $\C$. По умолчанию отождествляем комплесные числа с их координатными представлениями, а также функции с их координатными представлениями: $z = x + iy,\, f(z) = u(z) + i v(z) = u(x, y) + i v(x, y)$.

\begin{definition}[Голоморфная функция]
    Функция $f: \C \to \C$ называется \textit{голоморфной в области} $G \subset \C$, если $\forall z_0 \in G \,\exists \underset{z \to z_0}{\lim} \frac{f(z) - f(z_0)}{z - z_0}$.
\end{definition}

\begin{definition}[Комплексная производная]
    Если предел $\underset{z \to z_0}{\lim} \frac{f(z) - f(z_0)}{z - z_0}$ существует, он называется \textit{производной} функции $f$ в точке $z_0$ и обозначается $f'$.
\end{definition}

Из существования такого предела следует условие Коши-Римана: рассмотрим $\varepsilon \in \R$. Тогда:
\[
    \underset{\varepsilon \to 0}{\lim} \frac{f(z_0 + \varepsilon) - f(z_0)}{\varepsilon} = u_x'(x_0, y_0) + i v_x'(x_0, y_0)
\]
\[
    \underset{\varepsilon \to 0}{\lim} \frac{f(z_0 + i \varepsilon) - f(z_0)}{i \varepsilon} = v_y'(x_0, y_0) - i u_y'(x_0, y_0)
\]

Поскольку эти пределы должны быть равны, должны быть равны их вещестенные и комплексные компоненты. Таким образом выполняются равенства $u_x' = v_y', u_y' = -v_x'$, то есть в точности условия Коши-Римана.

Аналогично вещественному случаю, выполнена формула
\[
    f(z) = f(z_0) + f'(z_0)(z - z_0) + o(|z - z_0|)
\]
и её выполнение в точности эквивалентно дифференцируемости $f$ в точке $z_0$.

Верно и обратное: если для $u, v$ в точке $z_0$ выполнены условия Коши-Римана, то $f$ дифференцируема в $z_0$:
\[
    u_x' = v_y' = a, \quad u_y' = -v_x' = -b
\]
\[
    u(x_0 + \Delta x, y_0 + \Delta y) = u(x_0, y_0) + a(x_0, y_0) \Delta x - b(x_0, y_0) \Delta y + o(\sqrt{\Delta x^2 + \Delta y^2})
\]
\[
    v(x_0 + \Delta x, y_0 + \Delta y) = v(x_0, y_0) + b(x_0, y_0) \Delta x + a(x_0, y_0) \Delta y + o(\sqrt{\Delta x^2 + \Delta y^2})
\]
Домножив второе равенство на $i$ и прибавив результат к первому равенству, получим в точности формулу, эквивалентную дифференцируемости $f$ в $z_0$.


\begin{definition}[Аналитическая функция]
    Функция $f: \C \to \C$ называется \textit{аналитической в области} $G \subset \C$, если $\forall z_0 \in G \,\exists R > 0 \,\exists \{a_i\} \,\forall z \in B_R(z_0) \big(f(z) = \sum_{k = 0}^{\infty} a_k (z - z_0)^k\big)$
\end{definition}

\begin{theorem}
\label{analit_holom}
    $f: \C \to \C$, $G \subset \C$ - область. $f$ - аналитична в $G$ $\iff$ $f$ - голоморфна в $G$.
\end{theorem}
\begin{proof}
    $\Ra$: Если функция $f$ - аналитична, она локально предсталяется в виде степенного ряда. Поскольку этот ряд сходится равномерно, его можно почленно дифференцировать. Тогда у него существует производная, а значит, функция $f$ - голоморфна.
    
    $\La$: докажем позже.
\end{proof}

\begin{proposition}
    $f$ аналитична в $G$ $\Ra$ $\forall z_0 \in G \,\exists r > 0 \,\forall z, |z - z_0| < r \,\forall \partial D \subset B_r(z_0)$\\ $\left( \int_{\partial D} f(z)dz = 0 \right)$.
\end{proposition}
\begin{proof}
    $f$ аналитична $\Ra$ $f(z) = \sum a_n (z - z_0)^n$. Не умаляя общности, будем считать $z_0 = 0$.
    \begin{equation*}
        z^n dz = d\left( \frac{z^{n+1}}{n+1} \right) \Ra \int_{\partial D} z^n dz = \int_D d(z^n dz) = \int_D d\left(d\left(\frac{z^{n+1}}{n+1}\right)\right) = \int_D 0 = 0
    \end{equation*}
    По аддитивности интеграла, $\int_D f(z)dz = 0$.
\end{proof}

\begin{proposition}
\label{4eqiuv}
    $f$ определена в открытом круге. Тогда в этом круге следующие утверждения эквивалентны:
    \begin{enumerate}
        \item $\int_{\gamma} f(z)dz$ зависит только от $\gamma(0), \gamma(1)$.
        \item Если $\gamma$ - замкнутый путь, то $\int_{\gamma} f(z)dz = 0$.
        \item У $f(z)dz$ существует первообразная.
        \item Для любого прямоугольника $P$ со сторонами, параллельными осям, $\int_{\partial P} f(z)dz = 0$.
    \end{enumerate}
\end{proposition}
\begin{proof}
    $(1) \iff (2)$: очевидно. \\
    $(3) \Ra (2)$: очевидно. \\
    $(2) \Ra (4)$: очевидно. \\
    $(4) \Ra (3)$: предъявим кандидата на первообразную: $F(z) = \int_0^x f(s)ds + \int_{x}^{x+iy} f(s)ds$, где $z = x + iy$. Проверим, что производная $F$ равняется $f$. Для этого достаточно показать, что для $F$ выполняются условия Коши-Римана, и частные производные $F$ по $x, y$ равны $f(z)$:
    \begin{gather*}
        \underset{\varepsilon \to 0}{\lim} \frac{F(z + i\varepsilon) - F(z)}{i\varepsilon} = \underset{\varepsilon \to 0}{\lim} \frac{1}{i\varepsilon} \int_{[z, z + \varepsilon]} f(s)ds = \underset{\varepsilon \to 0}{\lim} \frac{1}{i\varepsilon} \int_0^1 f(z + ti\varepsilon)dt = \\
         = \frac{1}{i\varepsilon} \int_0^1 f(z + ti\varepsilon)dt = \frac{1}{i} \int_0^{\varepsilon} f(z + ti)dt = f(z)
    \end{gather*}
    где последнее равенство выполнено в силу равномерной непрерывности $f$ на отрезке. Теперь рассмотрим предел $\frac{F(z + \varepsilon) - F(z)}{\varepsilon}$. В силу условия (4)
    \begin{equation*}
        F(z) = \int_0^x f(s)ds + \int_{x}^{x+iy} f(s)ds = \int_{0}^{iy} f(s)ds + \int_{iy}^{x + iy} f(s)ds
    \end{equation*}

    \begin{center}
    \begin{tikzpicture}
        \begin{scope}[very thick, decoration={markings, mark=at position 0.5 with {\arrow{>}}}, >=Stealth]
            \draw[postaction={decorate}] (0,0) -- (4,0);
            \draw[postaction={decorate}] (4,0) -- (4,2);
            \draw[dashed] (4,2) -- (0,2);
            \draw[dashed] (0,2) -- (0,0);

            \draw[dashed] (6,0) -- (10,0);
            \draw[dashed] (10,0) -- (10,2);
            \draw[postaction={decorate}] (6,2) -- (10,2);
            \draw[postaction={decorate}] (6,0) node[below left]{$0$} -- (6,2);
        \end{scope}

        \fill (0,0) circle[radius=2pt] node[below left]{$0$};
        \fill (6,0) circle[radius=2pt] node[below left]{$0$};
        \fill (0,2) circle[radius=2pt] node[above left]{$iy$};
        \fill (6,2) circle[radius=2pt] node[above]{$iy$};
        \fill (4,0) circle[radius=2pt] node[below right]{$x$};
        \fill (10,0) circle[radius=2pt] node[below right]{$x$};
        \fill (4,2) circle[radius=2pt] node[above]{$x+iy$};
        \fill (10,2) circle[radius=2pt] node[above right]{$x+iy$};

        \draw [->,decorate,decoration={snake,amplitude=.4mm,segment length=2mm,post length=1mm}, thick]
            (4.5,1) -- (5.5,1);
    \end{tikzpicture}
    \end{center}

    Тогда, аналогично первому пределу, получаем, что
    \begin{gather*}
        \lim_{\varepsilon \to 0} \frac{F(z + \varepsilon) - F(z)}{\varepsilon} = \lim \frac{1}{\varepsilon} \int_{[z, z+\varepsilon]} f(z)dz = f(z)
    \end{gather*}
\end{proof}

\begin{remark}
    Будем говорить, что функция голоморфна/аналитична на множестве $M$, если найдётся область $G$ такая, что $M \subset G$ и функция голоморфна/аналитична в области $G$.
\end{remark}

\begin{theorem}
\label{circ_holom}
    $f$ голоморфна в замкнутом круге $\overline{B_r}$ $\Ra$ для любого прямоугольника $P \subset B_r$ со сторонами, параллельными осям, $\int_{\partial P} f(z)dz = 0$.
\end{theorem}
\begin{proof}
    Предположим противное, то есть нашёлся такой прямоугольник $P$, что $\int_{\partial P} f(z)dz = a \neq 0$. Разделим $P$ на четыре меньших прямоугольника:

    \begin{center}
    \begin{tikzpicture}
        \begin{scope}[very thick, decoration={markings, mark=at position 0.5 with {\arrow{>}}}, >=Stealth]
            \draw[postaction={decorate}] (0,0) -- node[below]{$B$} (4,0);
            \draw[postaction={decorate}] (4,0) -- (4,2);
            \draw[postaction={decorate}] (4,2) -- (0,2);
            \draw[postaction={decorate}] (0,2) -- node[left]{$A$} (0,0);
            \path (1.2,1.5) node{$P$};
            \path (10,2) node{$P_2$};
            \path (7,2) node{$P_1$};
            \path (7,0) node{$P_3$};
            \path (10,0) node{$P_4$};

            \path (7,1) node{$+$};
            \path (10,1) node{$+$};
            \path (8.5,0) node{$+$};
            \path (8.5,2) node{$+$};
            
            \foreach \y in {-0.5, 1.5}
                \foreach \x in {6, 9}
                {
                    \begin{scope}[xshift=\x cm, yshift=\y cm]
                        \draw[postaction={decorate}] (0,0) -- (2,0);
                        \draw[postaction={decorate}] (2,0) -- (2,1);
                        \draw[postaction={decorate}] (2,1) -- (0,1);
                        \draw[postaction={decorate}] (0,1) -- (0,0);
                    \end{scope}
                }
            
        \end{scope}

        \draw [->,decorate,decoration={snake,amplitude=.4mm,segment length=2mm,post length=1mm}, thick]
            (4.5,1) -- (5.5,1);
        \draw[dashed] (2,0) -- (2,2) (0,1) -- (4,1);
        \path (6,0) node[left]{$\frac{A}{2}$};
        \path (6,2) node[left]{$\frac{A}{2}$};
        \path (7,-0.5) node[below]{$\frac{B}{2}$};
        \path (10,-0.5) node[below]{$\frac{B}{2}$};
    \end{tikzpicture}
    \end{center}

    \begin{equation*}
        \int_{\partial P} f(z)dz = \underset{i}{\sum} \int_{\partial P_i} f(z)dz = a \Ra \exists i \left( \left| \int_{\partial P_i} f(z)dz \right| \geq \frac{|a|}{4} \right) 
    \end{equation*}

    Обозначим $P^1 = P_i$, после чего разделим уже $P_i$ на 4 части аналогичным образом. Часть с наибольшим модулем интеграла обозначим $P^2$. Продолжая прямоугольники с наибольшим модулем интеграла, получим последовательность вложенных прямоугольников $P^1, P^2, \hdots$ причём $\left| \int_{\partial P^n} f(z)dz \right| \geq \frac{|a|}{4^n}$. Отметим, что если стороны $P$ имели длину $A, B$, то стороны $P^n$ будут иметь длину $\frac{A}{2^n}, \frac{B}{2^n}$. Обозначим $\bigcap_n P^n = z_0$.

    \begin{gather*}
        f(z) = f(z_0) + f'(z_0)(z - z_0) + \varphi(z) \\
        \int_{\partial P^n} f(z)dz = \int_{\partial P^n} \Big( f(z_0) + f'(z_0)(z - z_0) \Big) + \int_{\partial P^n} \varphi(z)dz
    \end{gather*}

    Первый интеграл является интегралом многочлена по замкнутому контуру; у многочлена есть первообразная, а значит, по предложению \ref{4eqiuv}, этот интеграл равняется 0. Оценим второй интеграл: поскольку $\varphi = o(|z - z_0|)$, начиная с какого-то $n$ $|\varphi(z)| \leq \varepsilon |z - z_0|$. Тогда

    \begin{equation*}
        \left| \int_{\partial P^n} \varphi(z)dz \right| \leq \int_{\partial P^n} \left| \varphi(z)dz \right| \leq \frac{2A + 2B}{2^n} \cdot \varepsilon \frac{A+B}{2^n} \leq 2(A+B)^2 \frac{\varepsilon}{4^n}
    \end{equation*}

    Поскольку эта оценка верна для любого $\varepsilon$, устремляя $\varepsilon$ к 0, получаем $\int_{\partial P^n} \varphi(z)dz = 0$, а значит, и $\int_{\partial P^n} f(z)dz = 0$. Но по предположению $\left| \int_{\partial P^n} f(z)dz \right| \geq \frac{|a|}{4^n}$. Получили противоречие, значит, утверждение теоремы выполнено.

\end{proof}

\begin{lemma}[об устранимой особенности]
\label{noness}
    $z_0 \in \C$, $R > 0$. $f$ - непрерывна и ограничена в $B_R(z_0) \backslash \{z_0\}$. Для любого прямоугольника $P \subset B_R(z_0)$ со сторонами, параллельными осям, такого что $z_0 \not\in P$, $\int_{\partial P} f(z)dz = 0$. Тогда это же верно и для $P$ таких, что $z_0 \in P$.
\end{lemma}
\begin{remark}
    Если заменить шар $B_R(z_0)$ на его замыкание, лемма останется верной и доказывается аналогично.
\end{remark}
\begin{proof}
    Достаточно доказать для $P$ таких, что $z_0 \in \partial P$, так как если $z_0$ лежит строго внутри $P$, то $P$ можно разделить на два прямоугольника $P_1, P_2$ таких, что $z_0 \in \partial P_1, \partial P_2$:

    \begin{center}
    \begin{tikzpicture}
        \begin{scope}[very thick, decoration={markings, mark=at position 0.6 with {\arrow{>}}}, >=Stealth]
            \draw[postaction={decorate}] (0,0) -- (3,0);
            \draw[postaction={decorate}] (3,0) -- (3,2);
            \draw[postaction={decorate}] (3,2) -- (0,2);
            \draw[postaction={decorate}] (0,2) -- (0,0);
            \path (1.2,1.3) node{$P$};
            \draw[postaction={decorate}] (5,0) -- (7,0);
            \draw[postaction={decorate}] (7,0) -- (7,2);
            \draw[postaction={decorate}] (7,2) -- (5,2);
            \draw[postaction={decorate}] (5,2) -- (5,0);
            \path (6,1.3) node{$P_1$};
            \path (7.5,1) node{$+$};
            \draw[postaction={decorate}] (8,0) -- (9,0);
            \draw[postaction={decorate}] (9,0) -- (9,2);
            \draw[postaction={decorate}] (9,2) -- (8,2);
            \draw[postaction={decorate}] (8,2) -- (8,0);
            \path (8.5,1.3) node{$P_2$};
        \end{scope}
        \fill (2,0.5) circle[radius=2pt] node[below right]{$z_0$};
        \draw[dashed] (2,0) -- (2,2);

        \draw [->,decorate,decoration={snake,amplitude=.4mm,segment length=2mm,post length=1mm}, thick]
            (3.5,1) -- (4.5,1);

        \fill (7,0.5) circle[radius=2pt] node[left]{$z_0$};
        \fill (8,0.5) circle[radius=2pt] node[right]{$z_0$};
    \end{tikzpicture}
    \end{center}

    Заметим, что $P$ можно разделить на две части: тонкую полосу шириной $\delta$ $Q_{\delta}$, которая содержит на своей границе $z_0$, и остальная часть прямоугольника $P_{\delta}$. Поскольку $P_{\delta}$ не содержит $z_0$, по условию интеграл $f$ по $\partial P_{\delta}$ будет равен 0. Значит, для всякого $\delta$ интеграл по $\partial P$ равен интегралу по $\partial Q_{\delta}$.

    \begin{center}
    \begin{tikzpicture}
        \begin{scope}[very thick, decoration={markings, mark=at position 0.6 with {\arrow{>}}}, >=Stealth]
            \draw[postaction={decorate}] (0,0) -- (3,0);
            \draw[postaction={decorate}] (3,0) -- (3,2);
            \draw[postaction={decorate}] (3,2) -- (0,2);
            \draw[postaction={decorate}] (0,2) -- (0,0);
            \path (1.2,1.3) node{$P$};
            \draw[postaction={decorate}] (5,0) -- (7,0);
            \draw[postaction={decorate}] (7,0) -- (7,2);
            \draw[postaction={decorate}] (7,2) -- (5,2);
            \draw[postaction={decorate}] (5,2) -- (5,0);
            \path (6,1.3) node{$P_{\delta}$};
            \path (7.5,1) node{$+$};
            \draw[postaction={decorate}] (8,0) -- (9,0);
            \draw[postaction={decorate}] (9,0) -- (9,2);
            \draw[postaction={decorate}] (9,2) -- (8,2);
            \draw[postaction={decorate}] (8,2) -- (8,0);
            \path (9.5,1.5) node{$Q_{\delta}$};
        \end{scope}

        \fill (3, 0.5) circle[radius=2pt] node[below right]{$z_0$};
        \fill (9, 0.5) circle[radius=2pt] node[below right]{$z_0$};
        \draw [->,decorate,decoration={snake,amplitude=.4mm,segment length=2mm,post length=1mm}, thick]
            (3.5,1) -- (4.5,1);
        \draw[<->, semithick] (8,1.5) -- node[below]{$\delta$} (9,1.5);
    \end{tikzpicture}
    \end{center}

    Так как $P$ --- компакт, а $f$ непрерывна на $P$, то $f$ и равномерно непрерывна на $P$. Интеграл по $\partial Q_{\delta}$ равен сумме четырёх интегралов по сторонам. Две стороны имеют длину $\delta$, и по равномерной непрерывности $f$ модуль интеграла по ним оценивается сверху $2 \cdot \delta \varepsilon$, где $\varepsilon$ из определения равномерной непрерывности. Оставшиеся две стороны проходятся в разных направлениях и лежат на расстоянии $\delta$, а значит, модуль разности интегралов по ним тоже линейно оценивается через $\varepsilon$ в силу равномерной непрерывности:

    \begin{gather*}
        \left| \int_{\partial P} f(z)dz \right| = \left| \int_{\partial Q_{\delta}} f(z)dz \right| \leq 2 \delta \varepsilon + \left| \int_a^b f(t + \delta) - f(t) dt \right| \\
        \left| \int_a^b f(t + \delta) - f(t) dt \right| \leq \int_a^b |f(t + \delta) - f(t)| dt \leq (b - a) \varepsilon
    \end{gather*}

    Поскольку при $\delta \to 0$ $\varepsilon \to 0$, $\int_{\partial P} f(z)dz = 0$, что и требовалось доказать.
\end{proof}

\begin{proposition}
    $f$ голоморфна в круге $\Ra$ $\oint_{\gamma} f(z)dz = 0$.
\end{proposition}
\begin{proof}
    Достаточно применить теорему \ref{circ_holom} и предложение \ref{4eqiuv}.
\end{proof}

\begin{proposition}
    $f$ голоморфна в $G$, $a \in \overline{B_r(a)} \subset G$, $\gamma \subset B_r(a)$ - простой замкнутый путь. Тогда для любого $z$ внутри области, которую ограничивает $\gamma$, выполнено
    \begin{equation*}
        f(z) = \frac{1}{2 \pi i} \int_\gamma \frac{f(w)}{w - z} dw
    \end{equation*}
\end{proposition}
\begin{proof}
    $G$ --- открытая область, $\overline{B_r(a)} \subset G \Ra \exists \varepsilon > 0 \left( B_{r + \varepsilon}(a) \subset G \right)$.

    Зафиксируем $z$ внутри $\gamma$ и определим функцию $g: B_{r + \varepsilon}(a)\backslash\{z\} \to \C$, $g(w) = \frac{f(z) - f(w)}{z - w}$. $g$ голоморфна в $B_{r + \varepsilon}(a)\backslash\{z\}$, значит, по аналогии с теоремой \ref{circ_holom}, для любого прямоугольника $P \subset B_{r + \varepsilon}(a)\backslash\{z\}$ со сторонами, параллельными осям, $\int_{\partial P} g(w)dw = 0$. Значит, можно воспользоваться леммой \ref{noness} и получить, что это верно и для любого $P \subset B_{r + \varepsilon}(a)$. По предложению \ref{4eqiuv}, это эквивалентно тому, что $\int_{\gamma} g(w)dw = 0$ для любого замкнутого пути $\gamma \subset B_{r + \varepsilon}(a)$. Обозначим $S_r = \partial B_r(a) \subset B_{r + \varepsilon}(a)$ --- это простой замкнутый путь, обходимый против часовой стрелки. Тогда

    \begin{gather*}
        \int_{S_r} \frac{f(z)}{z - w} - \int_{S_r} \frac{f(w)}{z - w} = \int_{S_r} \frac{f(z) - f(w)}{z - w}dw = \int_{S_r} g(w)dw = 0 \Ra \\
        \Ra \int_{S_r} \frac{f(w)}{z - w}dw = \int_{S_r} \frac{f(z)}{z - w}dw = f(z) \int_{S_r} \frac{dw}{z - w}
    \end{gather*}

    Вычислим $\int_{S_r} \frac{dw}{z - w}$. Заметим, что $w \in S_r \iff |w - a| = r > |z - a|$. Тогда

    \begin{align*}
        \frac{1}{z - w} & = \frac{1}{(z - a) - (w - a)} \\&= - \frac{1}{w - a} \cdot \frac{1}{1 - \frac{z - a}{w - a}} \\&= - \frac{1}{w - a} \underset{n = 0}{\overset{\infty}{\sum}} \left( \frac{z - a}{w - a} \right)^n \\&= - \underset{n = 0}{\overset{\infty}{\sum}} \frac{(z - a)^n}{(w - a)^{n+1}}
    \end{align*}

    Посчитаем $\int_{S_r}$ от каждого слагаемого этого ряда. Для этого достаточно посчитать интеграл от знаменателя каждого слагаемого, так как числители не зависят от переменной интегрирования $w$.

    \begin{gather*}
        w = a + r e^{it},\, t \in [0, 2\pi) \Ra \int_{S_r} \frac{dw}{(w - a)^n} = \int_{0}^{2\pi} \frac{d\left(a + re^{it}\right)}{r^ne^{int}} = \int_{0}^{2\pi} \frac{d\left(re^{it}\right)}{r^ne^{int}} =\\= r^{-n+1} \int_{0}^{2\pi} e^{-int} d\left(e^{it}\right) = i r^{-n+1} \int_{0}^{2\pi} e^{-int} dt = \begin{cases}
            0, & n \neq 1 \\ 2\pi i, & n = 1
        \end{cases}
    \end{gather*}
    
    Таким образом получаем, что
    \begin{gather*}
        \int_{S_r} \frac{dw}{z - w} = - \underset{n = 0}{\overset{\infty}{\sum}} \int_{S_r} \frac{(z - a)^n}{(w - a)^{n+1}} dw = - 2\pi i \Ra \\
        \Ra \int_{S_r} \frac{f(w)}{z - w}dw = \int_{S_r} f(z) \frac{1}{z - w}dw = (- 2\pi i) f(z) \Ra\\\Ra f(z) = \frac{1}{2\pi i} \int_{S_r} \frac{f(w)}{w - z}dw
    \end{gather*}
\end{proof}

\begin{remark}
    \begin{align*}
        \frac{1}{w - z} = \underset{n = 0}{\overset{+\infty}{\sum}} \frac{(z - a)^n}{(w - a)^{n+1}} \Ra f(z) = \underset{n = 0}{\overset{+\infty}{\sum}} \frac{1}{2 \pi i} \left(\int_{S_r} \frac{f(w)}{(w - a)^{n+1}} dw\right) (z - a)^n = \underset{n = 0}{\overset{+\infty}{\sum}} c_n (z - a)^n
    \end{align*}
    Получили разложение $f$ в степенной ряд, значит, $f$ --- аналитическая, и теорема \ref{analit_holom} доказана.

    Более того, мы получили явное выражение для $c_n$, которым далее будем пользоваться:
    \begin{equation*}
        c_n = \frac{1}{2 \pi i} \left(\int_{S_r} \frac{f(w)}{(w - a)^{n+1}} dw\right)
    \end{equation*}
    Вообще говоря, эта формула выполняется и при $n < 0$ для функций $f$, которые можно разложить в ряд Лорана.
\end{remark}

\begin{remark}
    В предыдущих рассуждениях мы рассматривали точку $a \in G$ и круг некоторого радиуса $r$ вокруг неё, в котором $f$ раскладывается в ряд. Выясним, насколько большим можно взять этот радиус. Обозначим супремум таких радиусов за $R$.

    Ясно, что $R \geq \operatorname{dist}(a, \partial G)$. В самом деле, если это не так, то можно рассмотривать точки на границе шара $B_R(a)$ и раскладывать $f$ в ряд в этих точках. Эти ряды будут совпадать в области с изначальным рядом, а значит, должны просто полностью совпадать с ним. При этом окрестности, в которых сходятся эти ряды, в объединении содержат шар с центром в $a$ и радиусом, который строго больше $R$.

    \begin{center}
    \begin{tikzpicture}
        \draw[dashed, very thick] (0,0) circle[radius=3cm];
        \draw (-2.2,2.2) node[above left]{$G$};
        \fill (1,0) circle[radius=2pt] node[left]{$a$};
        \draw[dashed, blue] (1,0) circle[radius=1cm];
        % \draw[dashed, blue] (1,0) circle[radius=1.3cm];
        \draw[dashed, red]
            (1,1) circle[radius=0.5cm]
            (0,0) circle[radius=1.6cm]
            (1.7,-0.7) circle[radius=1cm]
            (1.7,0.7) circle[radius=0.8cm];
        \fill
            (1,1) circle[radius=1.5pt]
            (0,0) circle[radius=1.5pt]
            (1.7,-0.7) circle[radius=1.5pt]
            (1.7,0.7) circle[radius=1.5pt];
        \draw [->,decorate,decoration={snake,amplitude=.4mm,segment length=2mm,post length=1mm}, thick]
            (3.5,0) -- (4.5,0);
        
        \begin{scope}[shift={(8,0)}]
            \draw[dashed, very thick] (0,0) circle[radius=3cm];
        \draw (-2.2,2.2) node[above left]{$G$};
        \fill (1,0) circle[radius=2pt] node[left]{$a$};
        % \draw[dashed, blue] (1,0) circle[radius=1cm];
        \draw[dashed, blue] (1,0) circle[radius=1.3cm];
        \draw[dashed, red]
            (1,1) circle[radius=0.5cm]
            (0,0) circle[radius=1.6cm]
            (1.7,-0.7) circle[radius=1cm]
            (1.7,0.7) circle[radius=0.8cm];
        \fill
            (1,1) circle[radius=1.5pt]
            (0,0) circle[radius=1.5pt]
            (1.7,-0.7) circle[radius=1.5pt]
            (1.7,0.7) circle[radius=1.5pt];
        \end{scope}
    \end{tikzpicture}
    \end{center}

    С другой стороны, видно, что $R \leq \operatorname{dist}(a, \partial G)$. Если бы это было не так, область аналитичности $f$ была бы равна $B_R(a) \union G \supsetneq G$, ведь в любой точке $B_R(a)$ функция $f$ по определению $R$ раскладывается в ряд.
\end{remark}

\begin{theorem}[Мореры]
\label{Morera}
    $G$ - область. Следующие утверждения эквивалентны:
    \begin{enumerate}
        \item $f$ аналитична в $G$
        \item $f$ голоморфна в $G$
        \item $f(z)dz$ локально точна в $G$
    \end{enumerate}
\end{theorem}
\begin{proof}
    $(1) \iff (2)$: утверждение теоремы \ref{analit_holom}.\\
    $(2) \Ra (3)$: достаточно применить теорему \ref{circ_holom} и предложение \ref{4eqiuv}.\\
    $(3) \Ra (2)$: локальная точность $f(z)dz$ означает, что в окрестности любой точки из $G$ найдётся функция $F$ такая, что $dF = f(z)dz$. Зафиксируем $z \in G$.
    \begin{gather*}
        F_x dx + F_y dy = dF = f(z)dz = f dx + i f dy \\
        \begin{cases}
            F_x = f \\
            F_y = if
        \end{cases} \Ra F_y = i F_x \\
        F = u + iv \Ra u_y + iv_y = i(u_x + iv_x) = -v_x + iu_x \Ra \begin{cases}
            u_x = v_y \\
            u_y = - v_x
        \end{cases}
    \end{gather*}
    то есть для $F$ выполняются условия Коши-Римана в окрестности $z$. Как мы уже знаем, из этого следует, что $F$ в этой окрестности голоморфна, а значит и аналитична, а тогда аналитична и $f$, а тогда $f$ тоже голоморфна в окрестности $z$.
\end{proof}

\begin{definition}[Гармоническая функция]
    $f: \C \to \C$ называется \textit{гармонической}, если её вещественная часть $u$ гармонична как функция из $\R^2 \to \R$.
\end{definition}

\begin{remark}
    Если $f$ аналитична, то она и гармонична:
    \begin{equation*}
        f = u + iv,\,
        \begin{cases}
            u_x = v_y \\
            u_y = - v_x
        \end{cases} \Ra u_x^2 + v_y^2 = 0 \Ra \triangle u = u_x^2 = 0
    \end{equation*}
\end{remark}

\begin{remark}
    \begin{equation}\tag{$\star$}
        f(z) = \frac{1}{2\pi} \int_{0}^{2\pi} f(z + re^{it})dt = \frac{1}{2\pi i} \int_{S_r} \frac{f(w)}{w - z}dw
    \end{equation}
\end{remark}

\begin{proposition}[Принцип максимума]
\label{prmax}
    $G$ --- ограниченная область, $f$ аналитична в $\overline{G}$. Тогда $\underset{z \in \overline{G}}{\max} |f(z)| = \underset{z \in \partial G}{\max} |f(z)|$.
\end{proposition}
\begin{proof}
    Неравенство $\underset{z \in \overline{G}}{\max} |f(z)| \geq \underset{z \in \partial G}{\max} |f(z)|$ очевидно. $\underset{z \in \overline{G}}{\max} |f(z)|$ достигается, так как $\overline{G}$ замкнута и ограничена, следовательно, компактна. Пусть он достигается в точке $z_0 \notin \partial G$. Тогда $dist(z_0, \partial G) = r > 0$, и $(B_r(z_0) \subset G)$. Если $\underset{z \in \overline{G}}{\max} |f(z)| > \underset{z \in \partial G}{\max} |f(z)|$, то из непрерывности $f$ следует, что $f(z_0) > \int_0^{2\pi} f(z_0 + re^{it})dt$, что противоречит $(\star)$.
\end{proof}

\begin{remark}
    Если $\exists z_0 \in G \left(|f(z_0)| = \max_{\partial G}|f(z)|\right)$, то, поскольку выполнено $(\star)$, в некоторой окрестности $z_0$ выполнено $f(z) = f(z_0)$. Так как это выполнено для любой точки, в которой достигается максимум, множество таких точек открыто. По непрерывности $f$ это множество также и замкнуто. Поскольку область $G$ связна, это множество есть всё $G$.

    Это рассуждение называют \textit{строгим принципом максимума}.
\end{remark}

\begin{proposition}
    $G$ --- связная область, $f, g$ аналитичны в $G$, и в открытом $U \subset G$ $f \equiv g$. Тогда $f \equiv g$ во всём $G$.
\end{proposition}
\begin{proof}
    $f - g$ --- аналитическая функция, $f - g \equiv 0$ на $U$. Обозначим за $N \subset G$ множество нулей этой функции. Из условия утверждения следует, что $\Int N \neq \varnothing$. Покажем, что множество $\Cl(\Int N) \subset N$ и открытое, и замкнутое --- тогда из связности $G$ будет следовать, что $N = G$.

    Замкнутость тривиальна. Покажем открытость. У любой точки из $\Int \Int N = \Int N$ по определению есть окрестность, лежащая в $\Int N \subset \Cl(\Int N)$.
    
    У точки из $\Fr (\Int N)$ такая окрестность есть, так как в такой точке можно разложить $f - g$ в ряд: окрестность, в которой этот ряд сходится, пересекается с $\Int N$; это пересечение двух открытых множеств, значит, оно открыто; ряд на открытом подмножестве тождественно равен 0, значит, он просто тождественно равен 0; значит, вся окрестность, в которой этот ряд сходится, лежит в $\Int N \subset \Cl(\Int N)$.
\end{proof}

\begin{proposition}
    $f, g$ аналитичны в $G$, $\{\lambda_n\} \subset \C$, $\lambda_n \longrightarrow a \in G$, $f(\lambda_n) = g(\lambda_n)$. Тогда $f \equiv g$ в $G$.
\end{proposition}
\begin{proof}
    $f - g = \sum c_n (z - a)^n$. $(f-g)(\lambda_n) = 0 \Ra (f-g)(a) = 0 \Ra c_0 = 0$. Пусть $c_{n_0}$ - первый ненулевой коэффициент. $f(z) = (z - a)^{n_0} \sum_{k = 0}^{\infty} c_{n_0 + k} (z - a)^k = (z - a)^{n_0} \tilde{f}(z)$. $\tilde{f}(\lambda_n) = 0 \Ra \tilde{f}(a) = 0 \Ra c_{n_0} = 0$. Значит, все $c_n = 0$, значит, $f - g \equiv 0$.
\end{proof}

\begin{definition}[Целая функция]
    $f$ --- целая $\iff$ $f$ аналитична на $\C$.
\end{definition}

\begin{theorem}[Лиувилля]
\label{liuv}
    $F$ --- целая ограниченная функция $\Ra$ $F \equiv const$.
\end{theorem}
\begin{proof}
    $F(z) = \sum c_n z^n$ сходится $\forall z \in \C$. Как уже известно,
    \begin{equation*}
        c_n = \frac{1}{2 \pi i} \int_{S_r} \frac{F(w)}{w^{n+1}} dw
    \end{equation*}
    \begin{equation*}
        |c_n| \leq \underset{|z| = r}{\max} |F(z)| \cdot \frac{1}{2 \pi} \cdot 2 \pi r \cdot \frac{1}{r^{n+1}} = \frac{\underset{|z| = r}{\max} |F(z)|}{r^n} \leq \frac{c}{r^n} \underset{r \to +\infty}{\longrightarrow} 0
    \end{equation*}

    Значит, $c_n = 0$ для $n > 0$.
\end{proof}

\begin{remark}
    Обозначим, $M_F(r) = \underset{|z| = r}{\max} |F(z)|$. Из доказательства следует, что если $M_F(z) \leq c r^N$, то при $n > N$ $c_n = 0$, то есть $F$ - многочлен $\deg \leq N$.
\end{remark}

\begin{theorem}[Основная теорема алгебры]
    У многочлена $\deg = n \geq 1$ ровно $n$ корней с учётом кратности.
\end{theorem}
\begin{proof}
    Достаточно найти 1 корень.
    \begin{equation*}
        |p(z)| = |a_n z^n + \dots + a_0| = \left|a_n z^n \left(1 + \frac{a_{n-1}}{a_n} \frac{1}{z} + \dots \right)\right| \geq a_n |z|^n \left(1 - \left|\frac{a_{n-1}}{a_n}\right| \frac{1}{|z|} - \dots\right) \geq \frac{a_n}{2} |z|^n
    \end{equation*}
    для достаточно больших $|z|$.

    Если у $p$ нет корней, $\frac{1}{p(z)}$ --- голоморфная ограниченная функция в $\C$. Тогда по теореме \ref{liuv} $\frac{1}{p(z)} \equiv const \Ra p(z) \equiv const$, что невозможно, если $\deg(p) > 0$.
\end{proof}

\begin{remark}
    \todo{что-то про Римоновы поверхности}.
\end{remark}

\begin{definition}[Гомотопия]
    Непрерывная функция $H: [a, b] \times [0, 1] \to G \subset \C$ называется гомотопией.
\end{definition}

\begin{remark}
    Далее в курсе гомотопией будет называться гомотопия с фиксированными концами, то есть $H(a, t) \equiv const$, $H(b, t) \equiv const$. Два пути $\gamma_0, \gamma_1: [a, b] \to G$ называются \textit{гомотопными}, если $\gamma_0(a) = \gamma_1(a)$, $\gamma_0(b) = \gamma_1(b)$, и существует гомотопия $H$ такая, что $H(x, 0) = \gamma_0$, $H(x, 1) = \gamma_1$. Отношение гомотопности путей обозначается $\gamma_0 \sim \gamma_1$.

    \begin{center}
    \begin{tikzpicture}
        
    \end{tikzpicture}
    \end{center}
\end{remark}

\begin{theorem}
    $f$ голоморфна в $G$, $\gamma_0 \sim \gamma_1 \Ra \int_{\gamma_0} f(z)dz = \int_{\gamma_1} f(z)dz$.
\end{theorem}
\begin{proof}
    $H([a, b] \times [0, 1])$ --- компакт. Выберем у него конечное покрытие $B_1, \dots, B_N \subset G$. По лемме Лебега для достаточно малого $\varepsilon > 0$ образ при $H$ любого прямоугольника со сторонами меньше $\varepsilon$ целиком попадает в какой-нибудь $B_l$.

    Разобьём $[a, b] \times [0, 1]$ на такие прямоугольники $\Delta_{jk}$, $1 \leq j \leq m$, $1 \leq k \leq n$.

    \begin{center}
    \begin{tikzpicture}
        
    \end{tikzpicture}
    \end{center}

    $H(\Delta_{jk}) \subset B_l$, $f$ аналитична в $B_l$, значит
    \begin{equation*}
        \int_{H(\partial ([a, b] \times [0, 1]))} f(z)dz = \underset{j, k}{\sum} \int_{H(\partial \Delta_{jk})} f(z)dz = 0
    \end{equation*}
    \begin{gather*}
        \int_{H(\partial([a, b] \times [0, 1]))} f(z)dz = \int_{H([a, b] \times \{0\})} + \int_{H(\{b\} \times [0, 1])} + \int_{H([b, a] \times \{1\})} + \int_{H(\{a\} \times [1, 0])} = \\
        = \int_{\gamma_0} + \int_{\gamma_i(b)} - \int_{\gamma_1} - \int_{\gamma_i(a)} = \int_{\gamma_0} f(z)dz + 0 - \int_{\gamma_1} f(z)dz - 0 = 0
    \end{gather*}
    \begin{equation*}
        \int_{\gamma_0} f(z)dz = \int_{\gamma_1} f(z)dz
    \end{equation*}
\end{proof}

\begin{definition}[Односвязность]
    Область $G \subset \C$ называется односвязной, если любая петля в $G$ гомотопна тривиальной.
\end{definition}

\begin{theorem}
    $f$ аналитична в кольце $A_{r, R} = \{z \,|\, r < |z - a| < R\}$. Тогда $f(z) = \sum_{n \in \Z} c_n (z - a)^n$.
\end{theorem}

\begin{remark}
    В таком представлении $\sum_{n < 0}$ называется главной частью, а $\sum_{n \geq 0}$ --- аналитической частью. Аналитическая часть аналитична в $B_R(a)$, главная часть аналитична в $A_{r, R}$, поскольку она представляется как ряд от $1 - \frac{z}{a}$.
\end{remark}

\begin{proof}
    $\gamma_0 = \partial B_{R - \varepsilon}(a)$, $\gamma_1 = - \partial B_{r + \varepsilon}(a)$ (обход окружности в обратном направлении, то есть по часовой стрелке), $\gamma_+ = [a + r + \varepsilon, a + R - \varepsilon]$, $\gamma_- = - \gamma_+$.

    \begin{center}
    \begin{tikzpicture}[scale=1.5]
        \draw[dashed] (0,0) circle [radius=3cm];
        \draw[dashed] (0,0) circle [radius=1cm];
        \draw (-2.2,2.2) node[above left]{$A_{r,R}$};
        \draw[fill] (0,0) node[below right]{$a$} circle [radius=1pt];
        \draw[fill] (1.5,-1.5) node[below right]{$z$} circle [radius=1pt];
        \begin{scope}[arrows = {- Stealth[length=10pt]}]
            \draw (0,2.8) node[below]{$\gamma_0$} arc [start angle=90, end angle=270, radius=2.8cm];
            \draw (0,-2.8) arc [start angle=-90, end angle=90, radius=2.8cm];
            \draw (0,-1.2) arc [start angle=-90, end angle=-270, radius=1.2cm];
            \draw (0,1.2) node[above]{$\gamma_1$} arc [start angle=90, end angle=-90, radius=1.2cm];
        \end{scope}
        \draw[{Stealth[length=10pt,harpoon]}-{Stealth[length=10pt,harpoon]}] 
            (1.2,0) -- node[above right]{$\gamma_+$} node[below left]{$\gamma_-$} (2.8,0);
    \end{tikzpicture}
    \end{center}

    Обозначим $\gamma = \gamma_0 \union \gamma_1 \union \gamma_+ \union \gamma_-$. Видно, что $\gamma \sim 0$, $f$ аналитична, значит, $\int_{\gamma} f(z)dz = 0$. Тогда $\int_{\gamma} \frac{f(z) - f(w)}{z - w}dw = 0$. В $A_{r, R}$ выполнена формула Коши, потому что $\gamma \sim \partial B_{\varepsilon}(z)$.
    \begin{equation*}
        \int_{\gamma} \frac{f(w)}{w - z} dw = 0
    \end{equation*}
    \begin{equation*}
        \int_{\gamma} \frac{f(w)}{w - z} dw = \int_{\gamma_0} \frac{f(w)}{w - z}dw + \int_{\gamma_1} \frac{f(w)}{w - z}dw
    \end{equation*}
    \begin{equation*}
        |w - a| = R - \varepsilon \Ra \frac{1}{w - z} = \frac{1}{(w - a) - (z - a)} = \frac{1}{w - a} \cdot \frac{1}{1 - \frac{z - a}{w - a}} = \underset{n = 0}{\overset{\infty}{\sum}} \frac{(z - a)^n}{(w - a)^{n+1}}
    \end{equation*}
    \begin{equation*}
        |w - a| = r + \varepsilon \Ra \frac{1}{w - z} = \frac{1}{(w - a) - (z - a)} = \frac{1}{z - a} \cdot \frac{1}{\frac{w - a}{z - a} - 1} = - \underset{n = 0}{\overset{\infty}{\sum}} \frac{(w - a)^n}{(z - a)^{n+1}}
    \end{equation*}
    \begin{equation*}
        \int_{\gamma_0} \frac{f(w)}{w - z}dw + \int_{\gamma_1} \frac{f(w)}{w - z}dw = \underset{n = 0}{\overset{\infty}{\sum}} \frac{f(w)}{(w - a)^{n+1}}(z - a)^n + \underset{n = -\infty}{\overset{-1}{\sum}} f(w) (w - a)^{n+1} (z - a)^n
    \end{equation*}
\end{proof}

\begin{definition}[Особенность]
    $G = B_r(a) \backslash \{a\}$, $f$ голоморфна в $G$. Тогда говорится, что у $f$ в точке $a$ \textit{особенность}. Особенности бывают трёх типов:
    \begin{enumerate}
        \item Устранимая особенность: $f$ ограничена в окрестности $a$.
        \item Полюс: $|f(z)| \to +\infty$ при $z \to a$.
        \item Существенная особенность: $\lim_{z \to a} |f(z)|$ не существует.
    \end{enumerate}
\end{definition}

\begin{remark}
    Видно, что это исчерпывающая классификация. Приведём примеры функций, аналитичных в $B_1(0) \backslash \{0\}$, с особенностями каждого из типов:
    \begin{enumerate}
        \item $f(z) = z^3$. $|f(z)| = |z^3| = |z|^3 < 1$ в $B_1(0) \backslash \{0\}$.
        \item $f(z) = \frac{1}{z^2}$. $|f(z)| = \frac{1}{|z|^2} \to +\infty$ при $z \to 0$.
        \item $f(z) = e^{1/z}$. $|f(x + 0i)| = e^{1/x} \to +\infty$ при $x \to 0$. При этом $|f(0 + iy)| = |e^{-iy}| = 1$. Таким образом, $|f(z)|$ не имеет предела при $z \to 0$.
    \end{enumerate}
\end{remark}

\begin{proposition}
    $G = B_r(a) \backslash \{a\}$, $f$ аналитична в $G$. Тогда
    \begin{enumerate}
        \item в точке $a$ у $f$ устранимая особенность $\iff$ в разложении $f$ главная часть $\equiv 0$.
        \item в точке $a$ у $f$ полюс $\iff$ в главной части разложения $f$ количество ненулевых коэффициентов больше нуля и конечно.
        \item в точке $a$ у $f$ существенная особенность $\iff$ в главной части разложения $f$ ненулевых коэффициентов бесконечно много.
    \end{enumerate}
\end{proposition}
\begin{proof}
    \begin{enumerate}
        \item В точке $a$ у $f$ устранимая особенность $\iff$ (по лемме \ref{noness}) $f$ аналитична в окрестности $a$ $\iff$ $f = \sum c_n (z - a)^n$ $\iff$ главная часть $\equiv 0$.
        \item $\Ra$: Пусть в точке $a$ у $f$ полюс. Рассмотрим функцию $\frac{1}{f(z)}$, $f$ непрерывна и $|f(z)| \to +\infty$ при $z \to a$, значит, $|f|$ отделено 0 в некоторой окрестности $a$, значит $\frac{1}{f}$ определена и ограничена в этой окрестности. Тогда у $\frac{1}{f}$ в точке $a$ устранимая особенность, и тогда, по пункту $(1)$, \begin{gather*}
            \frac{1}{f(z)} = \underset{n=0}{\overset{+\infty}{\sum}} d_n (z - a)^n = c_{d^*} (z - a)^{d^*} + \underset{n > d^*}{\sum} d_n (z - a)^n \Ra \\ 
            \Ra \left| \frac{1}{f(z)} \right| = (c_{d^*} + \underline{o}(1)) |z - a|^{d^*} \Ra |f(z)| \leq C|z - a|^{-d^*}
        \end{gather*}
        Воспользуемся известной формулой для коэффициентов $c_n$ в разложении $f$ и оценим их по модулю:
        \begin{align*}
            |c_n| &\leq \left|\frac{1}{2 \pi i} \int_{S_{\varepsilon}} \frac{f(w)}{(w - a)^{n+1}} dw \right| \\
            &\leq \frac{1}{2 \pi} \left|\int_{0}^{2 \pi} f(a + \varepsilon e^{it}) (\varepsilon e^{it})^{-n-1} \cdot \varepsilon i e^{it} dt\right| \\
            &\leq \frac{1}{2 \pi} \int_{0}^{2 \pi} |f(a + \varepsilon e^{it})| \varepsilon^{-n} dt \\
            &\leq \frac{1}{2 \pi} \int_{0}^{2 \pi} C |\varepsilon e^{it}|^{-d^*} \varepsilon^{-n} dt \\
            &\leq \frac{C}{2 \pi} \int_{0}^{2 \pi} \varepsilon^{-d^* - n} dt \\
            &\leq C \varepsilon^{-d^* - n} \underset{\varepsilon \to 0}{\longrightarrow} 0 \text{ при } n < d^*
        \end{align*}

        Поскольку формула для $c_n$ верна при любом $\varepsilon$, получаем, что при $n < d^*$ $c_n = 0$. Значит, количество ненулевых коэффициентов в главной части разложения $f$ не больше чем $d^*$.

        $\La$: \begin{equation*}
            f(z) = \sum_{n = -N}^{-1} c_n (z - a)^n + \sum_{n = 0}^{+\infty} c_n (z - a)^n \Ra |f(z)| \underset{z \to a}{\longrightarrow} +\infty
        \end{equation*}
        \item Поскольку классификация особенностей исчерпывающая, и была доказана эквивалентность в пунктах $(1)$ и $(2)$, эквивалентность для пункта (3) выполнена автоматически.
    \end{enumerate}
\end{proof}

\begin{definition}[Кратность полюса]
    Если $f$ имеет в точке $a$ полюс, $N = \max\{n \,|\, c_{-n} \neq 0\}$ называется \textit{кратностью полюса}.
\end{definition}

\begin{proposition}
    Кратность полюса $f$ равна кратности нуля $\frac{1}{f}$.
\end{proposition}

\begin{definition}[Мероморфная функция]
    $f$ голоморфна в $G$. $f$ называется \textit{мероморфной} в $G$, если множество полюсов $f$ в $G$ состоит только из изолированных точек.
\end{definition}

\begin{definition}[Вычет]
    Если $f = \sum c_n (z - a)^n$, то \textit{вычетом} $f$ в точке $a$ называется $c_{-1}$. Вычет обозначается как $\Res_a f$.
\end{definition}

\begin{lemma}
    \todo{Про обход полюса и вычет}
    \begin{equation*}
        \int_{S_{\varepsilon}} f(w)dw = 2 \pi i \cdot \Res_a f
    \end{equation*}
\end{lemma}

\begin{theorem}[О вычетах]
\label{res}
    $f$ аналитична в $G \backslash \{a_1, \dots, a_n\}$. $\gamma$ --- замкнутый путь в $G$, причем для $\forall k (\ind_{a_k}(\gamma) = 1)$. Тогда
    \begin{equation*}
        \int_{\gamma} f(z)dz = 2 \pi i \underset{k = 1}{\overset{n}{\sum}} \Res_{a_k} f
    \end{equation*}
\end{theorem}
\begin{proof}
    \todo{Гомотопия + лемма}

    \begin{center}
    \begin{tikzpicture}[scale=2, transform shape]
        \draw (0,0) pic{atp};
    \end{tikzpicture}
    \end{center}
\end{proof}

\begin{remark}
    В качестве примера использования полученной теоремы вычислим $\int_\R \frac{\cos x}{1 + x^2} dx$.
    \begin{gather*}
        \int_\R \frac{\cos x}{1 + x^2} dx = \Re \left( \int_\R \frac{e^{ix}}{1 + x^2} dx \right) \\
        \int_\R \frac{e^{ix}}{1 + x^2} = \lim_{R \to +\infty} \int_{-R}^{R} \frac{e^{ix}}{1 + x^2}
    \end{gather*}
    Функция $\frac{e^{iz}}{1 + z^2}$ аналитична в $\C \backslash \{i, -i\}$. Для применения теоремы \ref{res} замкнём контур верхней полуокружностью радиуса $R$, обозначим эту дугу $\gamma_+$. Обозначим получившийся контур $\gamma_R$.

    % \begin{center}
    % \begin{tikzpicture}
    %     \draw (-3,0) node[below]{$-R$} (0,0) -- (3, 0) node[below]{$R$};
    %     \draw[-{Stealth}] (-3,0) -- (0,0);
    %     \draw (0,1) node[below right]{$i$} pic{cross};
    %     \draw[-{Stealth}] (3,0) arc[start angle=0, end angle=90, radius=3];
    % \end{tikzpicture}
    % \end{center}

    По теореме \ref{res}
    \begin{gather*}
        \frac{e^{iz}}{1 + z^2} = \frac{1}{z - i} \frac{e^{iz}}{z + i} \Ra c_{-1} = \frac{e^{-1}}{2i} \Ra \\
        \Ra \int_\gamma \frac{e^{iz}}{1 + z^2} dz = 2 \pi i \Res_i f = \frac{\pi}{e}
    \end{gather*}
    Заметим, что
    \begin{equation*}
        \left| \int_{\gamma_+} \frac{e^{iz}}{1 + z^2} dz \right| \leq \frac{1}{R^2 - 1} \cdot \pi R \underset{R \to +\infty}{\longrightarrow} 0
    \end{equation*}
    Поскольку для любого $R$ $\int_{\gamma_R} = \int_{-R}^{R} + \int_{\gamma_+}$, при устремлении $R \to +\infty$ получаем что $\int_\R = \frac{\pi}{e}$.
\end{remark}

\begin{remark}
    Вычислим ещё один интеграл: $\int_\R \frac{\sin x}{x} dx$. Из прошлых семестров мы уже знаем, что он равен $\pi$ --- получим этот результат методами комплексного анализа.

    $\frac{\sin z}{z} = \frac{e^{iz} - e^{-iz}}{2iz}$ --- целая функция. Хотим замкнуть контур, но в точке $z = 0$ функция $\frac{e^{iz}}{z}$ не суммируема, поэтому будем интегрировать по $[-R, -\varepsilon] \union [\varepsilon, R]$, устремляя $R \to +\infty, \varepsilon \to 0$. Такую пару отрезков можно замкнуть половиной кольца $A_{\varepsilon, R}$. Назовём полученный контур $\gamma_{\varepsilon, R}$, верхнюю полуокружность $c_R^+$, нижнюю $c_\varepsilon^+$.

    % \begin{center}
    % \begin{tikzpicture}
        
    % \end{tikzpicture}
    % \end{center}

    \begin{equation*}
        \int_\R \frac{\sin x}{x} dx = \lim_{\substack{\varepsilon \to 0 \\ R \to +\infty}} \int_{\varepsilon < |x| < R} \frac{\sin x}{x} dx = \lim_{\substack{\varepsilon \to 0 \\ R \to +\infty}} \frac{1}{i} \int_{\varepsilon < |x| < R} \frac{e^{ix}}{x} dx
    \end{equation*}
    \begin{equation*}
        \int_{\gamma_{\varepsilon, R}} \frac{e^{iz}}{z} dz = 0 = \int_{[\hdots] \union [\hdots]} + \int_{c_\varepsilon^+} + \int_{c_R^+}
    \end{equation*}
    Для вычисления искомого интеграла вычислим интегралы по полуокружностям. Рассмотрим полуокружность радиуса $R$. Заделим её на три части: $\arg \in [0, \frac{1}{\sqrt{R}}]$, $\arg \in (\frac{1}{\sqrt{R}}, \pi - \frac{1}{\sqrt{R}})$, $\arg \in [\pi - \frac{1}{\sqrt{R}}, \pi]$. В первой и третьей частях значение подынтегральной функции по модулю $\leq 1$, но сами части малы; вторая же часть большая, но значение подынтегральной функции на ней можно оценить лучше:
    \begin{gather*}
        \int_{c_R^+} \frac{e^{iz}}{z} dz = \int_{0}^{\pi} \frac{e^{i(Re^{it})}}{Re^{it}} d(Re^{it}) = \int_{0}^{\pi} e^{-R \sin(t)} \cdot e^{iR \cos(t)} dt = \int_{0}^{\frac{1}{\sqrt{R}}} + \int_{\frac{1}{\sqrt{R}}}^{\pi - \frac{1}{\sqrt{R}}} + \int_{\pi - \frac{1}{\sqrt{R}}}^{\pi}
    \end{gather*}
    В первом и третьем интегралах подынтегральная функция по модулю $<1$, а во втором второй множитель на модуль не влияет.
    \begin{equation*}
        \int_{0}^{\frac{1}{\sqrt{R}}} + \int_{\frac{1}{\sqrt{R}}}^{\pi - \frac{1}{\sqrt{R}}} + \int_{\pi - \frac{1}{\sqrt{R}}}^{\pi} \leq \frac{1}{\sqrt{R}} + \left(\pi - 2 \frac{1}{\sqrt{R}}\right) \cdot e^{-R \sin(t)} + \frac{1}{\sqrt{R}} \underset{R \to +\infty}{\longrightarrow} 0
    \end{equation*}

    Теперь рассмотрим интеграл по второй полуокружности, $c_\varepsilon^+$.
    \begin{equation*}
        \int_{c_\varepsilon^+} \frac{e^{iz}}{z} dz = \int_{c_\varepsilon^+} \frac{dz}{z} + \int_{c_\varepsilon^+} \frac{e^{iz} - 1}{z} dz
    \end{equation*}
    В окрестности 0 $\frac{e^{iz} - 1}{z}$ ограничена. Поскольку при $\varepsilon \to 0$ длина $c_\varepsilon^+ \to 0$, это означает, что второй интеграл в этой сумме $\to 0$ при $\varepsilon \to 0$.
    \begin{equation*}
        \int_{c_\varepsilon^+} \frac{dz}{z} = \int_{\pi}^{0} \frac{d(\varepsilon e^{it})}{\varepsilon e^{it}} = - \pi i
    \end{equation*}

    Таким образом
    \begin{gather*}
        i \int_\R \frac{\sin x}{x} dx = \lim_{\substack{\varepsilon \to 0 \\ R \to +\infty}} \int_{[\hdots] \union [\hdots]} \frac{e^{iz}}{z} dz = \pi i \Ra \\
        \Ra \int_\R \frac{\sin x}{x} dx = \pi
    \end{gather*}
\end{remark}

\begin{remark}
    Как мы знаем,
    \begin{equation*}
        f(z) = \frac{1}{2 \pi i} \int_{\partial D} \frac{f(w)}{w - z} dw = \begin{cases}
            0, w \notin D \\ f(z) \in \Int(D)
        \end{cases}
    \end{equation*}
    В глаза бросается факт, что случай $w \in \partial D$ остался нерассмотренным. Оказывается, что можно ввысти понятие \textit{полувычета}, и такой интеграл будет равен $\frac{f(z)}{2}$.
\end{remark}

\begin{remark}
    Пусть $f$ аналитична в проколотой окрестности $\dot{U}_a$ точки $a$, а в самой точке $a$ имеет существенную особенность. Следующие две теоремы рассматривают то, какие значения может принимать $f$ на $\dot{U}_a$.
\end{remark}

\begin{theorem}[Пикара]
    $f(\dot{U}_a)$ есть всё $\C$, возможно без одной точки.
\end{theorem}

\begin{theorem}[Пикара]
    Пусть $f$ не аналитична, а мероморфна. Тогда $f(\dot{U}_a)$ есть всё $\hat{\C} = \C \union \{\infty\}$, возможно без одной или двух точек.
\end{theorem}

\begin{remark}
    Эти теоремы сложные, поэтому даны без доказательства. Есть более слабая, но и намного более простая теорема схожего характера.
\end{remark}

\begin{theorem}[Сохоцкого]
    $f$ мероморфна в $\dot{U}_a$, и имеет существенную особенность в $a$. Тогда $f(\dot{U}_a)$ всюду плотно в $\C$.
\end{theorem}
\begin{proof}
    Предположим противное, то есть $\exists B_r(\xi) \subset \C \backslash f(\dot{U}_a)$. Рассмотрим функцию $h(z) = \frac{1}{f(z) - \xi}$. Она ограничена и аналитична в $\dot{U}_a$, а значит, по лемме \ref{noness}, аналитична в $U_a$. Однако $f(z) = \frac{1}{h(z)} + \xi$, и если $h(z)$ аналитична в $U_a$, то $f(a)$ можно однозначно определить (если $h(a) = 0$, то значение получится равным $\infty$). Это противоречит тому, что в точке $a$ у $f$ была существенная особенность.
\end{proof}

\begin{remark}
    Пусть $f$ аналитична в области $G$. Попытаемся определить $\log(f(z))$ в $G$ --- вообще говоря, $\log$ функция не однозначная, потому что $\exp(\log(f(z)) + 2 i \pi k) = \exp(\log(f(z)))$. Тем не менее, можно определить $\log$ в области через фиксированное значение в точке и непрерывность. Для этого потребуем, что $f(z) \neq 0$ в $G$, и что $G$ односвязна. Если бы мы смогли задать $\log(f(z))$, то его производной должна была бы быть функция $\frac{f'(z)}{f(z)}$. Исходя из этого, зафиксируем точку $z_0 \in G$, и определим
    \begin{equation*}
        H(z) = \int_{z_0}^{z} \frac{f'(z)}{f(z)} dz
    \end{equation*}
    $H(z)$ определена корректно, ведь поскольку $f$ аналитична и $G$ односвязна, интегралы по разным путям из $z_0$ в $z$ совпадают. Покажем, что $H(z)$ почти является логарифмом:
    \begin{equation*}
        \left( \frac{e^{H(z)}}{f(z)} \right)' = \frac{e^{H(z)} H'(z) f(z) - e^{H(z)} f'(z)}{f(z)^2} = \frac{e^{H(z)} \frac{f'(z)}{f(z)} f(z) - e^{H(z)} f'(z)}{f(z)^2} = 0 \Ra e^{H(z)} = const \cdot f(z)
    \end{equation*}
    Заменив $H(z)$ на $H(z) + c$, можно получить $e^{H(z)} = f(z)$. Заметим, что любая другая функция с таким свойством отличается прибавленим $2 i \pi k$:
    \begin{equation*}
        \begin{cases}
            e^{H_1(z)} = f(z) \\ e^{H_2(z)} = f(z)
        \end{cases} \Ra e^{H_1(z) - H_2(z)} = 2 i \pi k
    \end{equation*}
    При этом $k$ не может меняться для фиксиованной пары функций $(H_1, H_2)$, потому что равенство выполнено при любом $z$ и $H_1 - H_2$ --- непрерывная функция.
\end{remark}

\begin{remark}
    Определив логарифм, можно определить и возведение комплексного числа в комплексную степень: $z^\beta = e^{\beta \log(z)}$.
\end{remark}

\begin{remark}
    Подумаем о "<физическом смысле"> $\int_{z_0}^{z} H(w) dw = \log(z) - \log(z_0)$. Пусть $f$ аналитична в $G$, $\{z_k\}$ --- нули $f$, $\gamma \subset G$ - путь не проходящий через $z_k$, и $\forall k (\ind_{z_k}(\gamma) = 1)$.

    Рассмотрим интеграл $\int_\gamma H(w) dw$. Отметим две близкие точки $a, b$ на контуре $\gamma$ следующим образом:

    % \begin{center}
    % \begin{tikzpicture}
        
    % \end{tikzpicture}
    % \end{center}

    Обозначим часть $\gamma$ от $a$ до $b$ ("<от"> и "<до"> понимаются в смысле обхода $\gamma$ в правильном направлении, то есть эта часть будет почти всем контуром) как $\gamma_{ab}$. Достаточно малая окрестность $\gamma_{ab}$ не содержит нулей $f$ и односвязна, значит в ней можно определить $\log(f(z))$ как интеграл $\int_{a}^{z} H(w) dw$. Интеграл по части контура $\gamma_{ab}$ равен $\log(f(b)) - \log(f(a))$. 
    
    Отметим, что если $f(z) = r e^{i \varphi}$, то $\log(f(z)) = \log(r) + i \varphi$. Заметим, что после обхода $\gamma$ $\arg f(z) = \varphi$ мог измениться на $2 i \pi k$. Обозначим это изменение $\Delta_{\arg f(z)}$.   
    
    Устремив точку $b$ к $a$, получим, что $\int_\gamma H(w) dw = i \Delta_{\arg f(z)}$.

    Заметим, что $\gamma$ гомотопно объединению обходов нулей:

    % \begin{center}
    % \begin{tikzpicture}
        
    % \end{tikzpicture}
    % \end{center}

    Теперь распишем интеграл $\int_\gamma H(w) dw = \int_\gamma \frac{f'(w)}{f(w)} dw$ на окружности малого радиуса вокруг $z_k$.
    \begin{gather*}
        f(z) = \sum_{l = 0}^{+\infty} c_l (z - z_k)^l = c_{l^*} (z - z_k)^{l^*} + \sum_{l > l^*} c_l (z - z_k)^l, \quad f(z_k) = 0 \Ra l^* \geq 1  \\
        f'(z) = l^* c_{l^*} (z - z_k)^{l^* - 1} + \sum_{l > l^*} l c_l (z - z_k)^{l - 1} \\
        \frac{f'(z)}{f(z)} = \frac{l^* c_{l^*} (z - z_k)^{l^* - 1} + (z - z_k)^{l^*} \left(\sum_{l > l^*} l c_l (z - z_k)^{l - 1 - l^*} \right)}{c_{l^*} (z - z_k)^{l^*} \left(1 + (z - z_k) \sum_{l > l^*} c_l c_{l^*}^{-1} (z - z_k)^{l - l^* - 1} \right)} = \\
        = \frac{l^* c_{l^*} + (z - z_k) \varphi(z)}{c_{l^*} (z - z_k) (1 + (z - z_k)\psi(z))} = \frac{l^*}{(z - z_k) (1 + (z - z_k)\psi(z))} + \frac{\varphi(z)}{c_{l^*} (1 + (z - z_k)\psi(z))}
    \end{gather*}
    Вторая дробь есть частное двух аналитических фунций, причём знаменатель не обращается в 0 по непрерывности (т.к. $z$  лежит на окружности малого радиуса вокруг $z_k$). Значит, вторая дробь --- аналитическая функция.
    \begin{gather*}
        \frac{l^*}{(z - z_k) (1 + (z - z_k) \psi(z))} = \frac{l^*}{z - z_k} + \left( \frac{l^*}{(z - z_k)(1 + (z - z_k) \psi(z))} - \frac{l^*}{z - z_k} \right) = \\
        = \frac{l^*}{z - z_k} + \frac{l^*(1 - 1 - (z - z_k) \psi(z))}{(z - z_k) (1 + (z - z_k)\psi(z))} = \frac{l^*}{z - z_k} + \frac{l^* \psi(z)}{1 + (z - z_k) \psi(z)}
    \end{gather*}
    Здесь, по аналогичным причинам, вторая дробь тоже является аналитической фунцией.

    Итого, получаем
    \begin{equation*}
        \frac{f'(z)}{f(z)} = \frac{l^*}{z - z_k} + g(z)
    \end{equation*}
    где $g(z)$ аналитична на рассматриваемой окружности. Теперь можно вычислить
    \begin{equation*}
        \int_{\partial B_\varepsilon(z_k)} \frac{f'(w)}{f(w)} dw = \int_{\partial B_\varepsilon(z_k)} \frac{l^*}{z - z_k} + \int_{\partial B_\varepsilon(z_k)} g(w) dw = 2 i \pi \cdot l^* + 0 = 2 i \pi \cdot l^*
    \end{equation*}
    При этом $l^*$ есть в точности кратность корня $z_k$ у $f$. Тогда получаем
    \begin{equation*}
        \int_\gamma \frac{f'(w)}{f(w)} dw = \sum_k \int_{\partial B_\varepsilon(z_k)} \frac{f'(w)}{f(w)} dw = \sum_k 2 i \pi \cdot l_k^* = 2 i \pi \cdot n(f)
    \end{equation*}
    где $n(f)$ --- количество нулей $f$ с учётом кратности.

    Используя оба полученных равенства, получаем утверждение:
    \begin{equation*}
        2 i \pi n(f) = \int_\gamma \frac{f'(w)}{f(w)} dw = i \Delta_{\arg f(z)}
    \end{equation*}
\end{remark}

\begin{theorem}[Принцип аргумента]
\label{prarg}
    $\Delta_{\arg f(z)} = 2 \pi n(f)$.
\end{theorem}

\begin{remark}
    Если считать $f$ не аналитической, а мероморфной, при этом $f(z) \neq 0, \infty$ на $\gamma$, то аналогично доказывается аналог этой теоремы для ммероморфных функций.
\end{remark}

\begin{theorem}[Принцип аргумента для мероморфной функции]
    $\Delta_{\arg f(z)} = 2 \pi (n(f) - p(f))$, где $p(f)$ --- количество полюсов $f$ внутри $\gamma$.
\end{theorem}

\begin{remark}
    Приведём несколько примеров использования принципа аргумента.

    Пусть $P(z)$ --- полином степени $n$. Тогда при обходе по достаточно большой окружности $\Delta_{\arg P(z)} = 2 \pi n$. Значит, в достаточно большом круге у $P(z)$ ровно $\frac{2 \pi n}{2 \pi} = n$ нулей, что даёт основную теорему алгебры.

    Пусть $f(z) = 2 z^3 + \frac{\sin z}{10} - 1 = 0$ в $B_1(0)$. Если $|z| < 1$, то $|\sin(z)| < e$. Тогда при обходе $\gamma = \partial B_{1 - \varepsilon}(0)$
    \begin{equation*}
        |2 z^3| > \left| \frac{\sin(z)}{10} - 1 \right| \Ra \Delta_{\arg f(z)} = 2 \pi \cdot 3
    \end{equation*}
    значит, $f(z)$ имеет ровно 3 корня в $B_1(0)$.
\end{remark}

\begin{theorem}[Руше]
\label{rouche}
    $G$ --- область, $\gamma \subset G$ --- простой путь, обходимый против часовой стрелки, $f, g$ --- аналитические в $G$, $|f(\xi)| > |g(\xi)|$ при $\xi \in gamma$. Тогда у $f$ и $f + g$ одинаковое число нулей внутри $\gamma$.
\end{theorem}
\begin{proof}
    Обозначим за $N_f, N_{f + g}$ количество нулей внутри $\gamma$ у функций $f, f + g$ сооответственно. Из доказательства принципа аргумента знаем, что
    \begin{gather*}
        N_f = \frac{1}{2 i \pi} \int_\gamma \frac{f'(w)}{f(w)} dw, \quad N_{f + g} = \frac{1}{2 i \pi} \int_\gamma \frac{f'(w) + g'(w)}{f(w) + g(w)} dw \\
        N_f - N_{f + g} = \frac{1}{2 i \pi} \int_\gamma \frac{f'g - g'f}{f (f + g)} dw = \frac{1}{2 i \pi} \int_\gamma \frac{\left(1 + \frac{g}{f}\right)'}{1 + \frac{g}{f}} dw = N_{1 + \frac{g}{f}}
    \end{gather*}
    Поскольку $|g(\xi)| < |f(\xi)|$, $\left|\frac{g}{f}\right| < 1$. Тогда все значения $\left(1 + \frac{g}{f}\right)(\xi)$ лежат справа от точки 0. Значит, $\ind_0 \left(1 + \frac{g}{f}\right) = 0 \Ra N_{1 + \frac{g}{f}} = \frac{1}{2 \pi} \Delta_{\arg f(z)} = 0$. 
\end{proof}

\begin{remark}
    Приведём пример использования этой теоремы. Пусть $f(z) = z^3, g(z) = \frac{\sin z}{10} - \frac{1}{3}$, $\gamma = \partial B_1(0)$. Тогда $|f(\xi)| = 1 > \frac{e}{10} + \frac{1}{3} \geq |g(\xi)|$. При этом у $f$, очевидно, один корень в $B_1(0)$. Значит, и у функции $z^3 + \frac{\sin z}{10} - \frac{1}{3}$ ровно один корень в $B_1(0)$.
\end{remark}

\begin{definition}
    $\hat{\C} = \C \union \{\infty\}$.
\end{definition}

\begin{remark}
    Пусть $f$ мероморфна, $f: \hat{\C} \to \hat{\C}$. \todo{Хотим придать смысл} $f(\infty)$.
    \begin{equation*}
        f(z) = \sum_{n = -\infty}^{+\infty} \frac{c_n}{z^n} = \sum_{n \geq 1}^{+\infty} + \sum_{n = -\infty}^{0} = I + \sum_{n = -\infty}^{0}
    \end{equation*}
    
    Если I нет, то $\lim_{z \to \infty} f(z) = c_0$.

    Если количество слагаемых в I конечно, то $\lim_{z \to \infty} = \infty$.

    Если же количество слагаемых в I бесконечно, то у $f$ в точке $\infty$ существенная особенность.

    Сфера Римана.
\end{remark}

\begin{remark}
    Пусть $f$ аналитична в $\hat{B}_R = \{ z \,|\, |z| > R\}$, т.е. в окрестности $\infty$. Сделаем в форме $f(z)dz$ замену переменной:
    \begin{gather*}
        w = \frac{1}{z} \Ra f\left(\frac{1}{w}\right) d\left(\frac{1}{w}\right) = - \frac{1}{w^2} f\left(\frac{1}{w}\right) dw \Ra \\ \Ra f(z) dz = \sum c_n z^n dz = - \sum c_n w^{-n-2} dw = \frac{- c_{-1}}{w} + (\hdots)
    \end{gather*}
    Таким образом, естественно считать $\Res_\infty f = - c_1$.
\end{remark}

\begin{theorem}
    $f$ аналитична в $\hat{\C} \backslash \{a_1, \hdots, a_n\}$, где $a_i \in \hat{\C}$. Тогда $\sum_{i = 1}^{n} \Res_{a_k} f = 0$.
\end{theorem}
\begin{proof}
    Рассморим контур $\gamma$, который обходит все особенности в $\C$ против часовой стрелки по одному разу. Сначала рассмотрим случай, когда $\infty$ --- не особая точка. Тогда
    \begin{equation*}
        2 i \pi \sum \Res_{a_k} f = \int_\gamma f(w) dw
    \end{equation*}
    Этот интеграл равен 0, так как $\gamma$ можно стянуть в точку через $\infty$. Значит, и сумма вычетов равна 0.

    Если же $\infty$ является особой точкой, то $\gamma$ гомотопен обходу $\infty$ по часовой стрелке. Если обозначить за $\{b_i\}$ все $a_i$ кроме $\infty$, получим $\sum \Res_{b_k} f = c_{-1} = -\Res_\infty f$. Тогда $\sum \Res_{b_k} f + \Res_\infty f = \sum \Res_{a_k} f = 0$.
\end{proof}

\begin{remark}
    Пусть $f$ аналитична в $G$. Рассмотрим её поведение в малой окрестности $U_a$ точки $a \in G$. Мы уже обсуждали несколько теорем о случае, когда $a$ --- существенная особенность. Теперь же обратимся к случаю, когда в точке $a$ нет особенностей.
    \begin{equation*}
        f(z) = f(a) + f'(a)(z - a) + g(z) \Ra f(z) - f(a) = f'(a)(z - a) + o(|z - a|), \quad g(z) = o(|z - a|)
    \end{equation*}
    Если $f'(a) \neq 0$, то в правой части первое слагаемое в малой окрестости $a$ по модулю больше второго, значит, по теореме \ref{rouche} для пары $(f'(a)(z - a), g(z))$ и $\gamma = \partial U_a$, у функции $f(z) - f(a)$ столько же нулей в $U_a$, сколько и у $f'(a)(z - a)$, то есть ровно 1. Таким образом, $f$ принимает значение $f(a)$ единожды в $U_a$. Это подводит нас к понятию однолистных функций.
\end{remark}

\begin{definition}[Однолистная функция]
    Пусть $G$ --- обсласть, $f$ аналитична в $G$. $f$ называется \textit{однолистной} в $G$, если она биективна, то есть $\forall z_1, z_2 \in G (f(z_1) = f(z_2) \Ra z_1 = z_2)$.
\end{definition}

\begin{remark}
    \todo{Преобразования Мёбиуса - автоморфизмы диска}
\end{remark}

% переделать выше: ref |-> название (ref)

\begin{lemma}[Шварца]
\label{lmsch}
    $D = B_1(0)$, $f: D \to D$ аналитична в $D$, $f(0) = 0$. Тогда $|f(z)| \leq |z|$, и если найдётся $z_0$ такое, что $|f(z_0)| = |z_0|$, то $f$ --- поворот.
\end{lemma}
\begin{proof}
    $f(z) = \sum_{n = 0}^{+\infty} c_n z^n$, $f(0) = c_0 = 0$, значит, $g(z) = \frac{f(z)}{z}$ аналитична в $D$. Рассмотрим $r < 1$. $|f(z)| < 1$, значит, при $|z| = r$ $|g(z)| < \frac{1}{r}$. По принципу максимума (\ref{prmax}) для $G = B_r(0)$ из этого следует что $|g(z)| < \frac{1}{r}$ и при $|z| \leq r$. Тогда $\forall z \in D (|g(z)| \leq 1)$. По определению $g(z)$ из этого следует $|f(z)| \leq |z|$.

    Если нашлось $z_0$ такое, что $|f(z_0)| = |z_0|$, то $|g(z_0)| = 1$. Тогда по строгому принципу максимума $|g(z)| = 1$ на всём $D$. 
\end{proof}

\begin{proposition}
    Пусть $f: D \to D$ --- аналитический автоморфизм $D$, и $f(0) = 0$. Тогда $f$ --- поворот.
\end{proposition}
\begin{proof}
    $f$ --- автоморфизм, следовательно, и $f^{-1}$ --- автоморфизм. Также $f^{-1}(0) = 0$. Рассмотрим $z_0 \in D,\, z_0 \neq 0$. По лемме Шварца (\ref{lmsch}), $|z_0| \geq |f(z_0)|$. Но лемму Шварца можно применить и к $f^{-1}$, из чего вытекает $|f(z_0)| \geq |f^{-1}(f(z_0))| = |z_0|$. Таким образом $|z_0| \geq |f(z_0)| \geq |z_0|$, а значит, достигается равенство $|f(z_0)| = |z_0|$. Тогда, опять же по лемме Шварца, $f$ --- поворот.
\end{proof}

\begin{remark}
    \todo{Автоморфизмы диска - преобразования Мёбиуса}
\end{remark}

\begin{definition}
    $A(G)$ --- векторное пространство функций, аналитичных в $G$. \\
    $\{f_n\}$ --- последовательность аналитических функций. Тогда $f_n \to f$ в $A(G)$ по поределению означает, что в любом компакте $K \subset G$ выполнено $f_n \rightrightarrows f$ на $K$.
\end{definition}

\begin{proposition}[Uniform convergence on compact sets, u.c.c.s]
\label{uccs}
    $\{f_n\}$ --- последовательность аналитических в $G$ функций, $f_n \to f$ в $A(G)$. Тогда $f$ --- аналитическая в $G$.
\end{proposition}
\begin{proof}
    Если для любого компакта $K \subset G$ $f \in C(K)$, то и $f \in C(G)$. Теперь для проверки аналитичности достаточно показать, что интегралы по замкнутым контурами $\gamma \subset G$ равны 0. Рассмотрим такой контур $\gamma$. Найдётся компакт $K \supset \gamma$ (например, подойдёт замыкание области, которую ограничивает $\gamma$). Поскольку все $f_n$ аналитичны на $K$, $\int_\gamma f_n(z)dz = 0$.
    \begin{equation*}
        \left| \int_\gamma f(z)dz - \int_\gamma f_n(z)dz \right| \leq \int_\gamma |f(z) - f_n(z)| |dz| \leq \underset{z \in K}{\sup} |f_n(z) - f(z)| \cdot L(\gamma) \underset{n \to +\infty}{\longrightarrow} 0
    \end{equation*}
    где предел равен 0 потому, что на $K$ $f_n \rightrightarrows f$.
\end{proof}

\begin{remark}
    Как уже было сказано, $A(G)$ является векторным пространством. Вообще говоря, на нём можно задать и топологию --- у нас есть определение сходимости, и замкнутыми множествами естественно назвать те, которые содержат свои предельные точки; соответственно, открытыми множествами будут их дополнения.
\end{remark}

\begin{theorem}
    $f_n \to f$ в $A(G)$ $\Ra$ $f_n' \to f'$ в $A(G)$.
\end{theorem}
\begin{proof}
    Нужно доказать, что на любом компакте $K \subset$ $f_n \rightrightarrows f$. Сначала докажем это не для $K$, а для $B_r \subset \overline{B_{r + \varepsilon}} \subset G$, то есть таких шаров, которые внутри $G$ и отделяются от него замкнутым шаром чуть большего радиуса.

    Зафиксируем $z \in B_r$. Обозначим $\gamma_\varepsilon = \partial B_\varepsilon(z)$. Отметим, что $\gamma_\varepsilon \subset \overline{B_{r + \varepsilon}} \subset G$. Тогда
    \begin{equation*}
        f_n'(z) - f'(z) = \frac{1}{2 i \pi} \int_\gamma \frac{(f_n - f)(w) dw}{(w - z)^2} = \frac{1}{2 i \pi} \int_\gamma \frac{(f_n - f) dw}{\varepsilon^2} \Ra |(f_n - f)'(z)| \leq \frac{\max_{z \in \gamma_\varepsilon} |(f_n - f)(z)|}{2 \pi \varepsilon^2}
    \end{equation*}
    \begin{equation*}
        \sup_{z \in B_r} |f_n'(z) - f'(z)| = \sup_{z \in B_r} |(f_n - f)'(z)| \leq \frac{\max_{z \in \gamma_\varepsilon} |(f_n - f)(z)|}{2 \pi \varepsilon^2} \rightrightarrows 0
    \end{equation*}
    где последний предел равен 0, так как $f_n \to f$ в $A(G)$.
\end{proof}

\begin{remark}
    Тогда если $f_n \to f$ в $A(G)$, то для любого $k \in \N$ $f_n^{(k)} \to f^{(k)}$ в $A(G)$.
\end{remark}

\begin{proposition}
    Условие $f_n \rightrightarrows f$ на $K$ эквивалентно условию $\int_K |f_n(z) - f(z)| d \lambda_2^*(z) \to 0$.
\end{proposition}
\begin{proof}
    $\Ra$: $\left| \int_K |f_n(z) - f(z)| d \lambda_2^*(z) \right| \leq \lambda_2^*(K) \cdot \max_{z \in K} |f_n(z) - f(z)| \to 0$.

    $\La$: рассмотрим $a \in K$. Обозначим $B_\varepsilon = B_\varepsilon(a)$, $K_\varepsilon$ --- $\varepsilon$-окрестность $K$. Как известно, $K \subset G \Ra \exists \varepsilon (K_\varepsilon \subset G)$. Тогда
    \begin{equation*}
        |f_n(a) - f(a)| = \left| \frac{1}{\lambda_2^*(B_\varepsilon)} \int_{B_\varepsilon} (f_n(w) - f(w)) d \lambda_2^*(w) \right| \leq \frac{1}{\lambda_2^*(B_\varepsilon)} \int_{B_\varepsilon} |f_n(w) - f(w)| d \lambda_2^*(w)
    \end{equation*}
    Поскольку $a \in K$, $B_\varepsilon \subset K_\varepsilon$. Тогда
    \begin{equation*}
        \frac{1}{\lambda_2^*(B_\varepsilon)} \int_{B_\varepsilon} |f_n(w) - f(w)| d \lambda_2^*(w) \leq \frac{1}{\lambda_2^*(B_\varepsilon)} \int_{K_\varepsilon} |f_n(w) - f(w)| d \lambda_2^*(w) \underset{n \to +\infty}{\longrightarrow} 0
    \end{equation*}
    Эта оценка не зависит от $a$, значит, $f_n \rightrightarrows f$ на $K$.
\end{proof}

\begin{remark}
    Если $\{f_n\}$ --- однолистные функции, и $f_n \to f$ в $A(G)$, то $f$ может не быть однолистной: например, если $f_n = c + \frac{1}{n z}$, то $f = c$. Однако оказывается, что это единственное возможно препятствие.
\end{remark}

\begin{theorem}[Об однолистности]
\label{one_leaf}
    $G$ --- свзяная область, $\{f_n\}$ --- однолистные в $G$ функции, $f_n \to f$ в $A(G)$. Тогда $f$ либо тождественно равна константе, либо однолистна.
\end{theorem}
\begin{proof}
    Предположим противное, то есть что $f$ не константа и не однолистна. Тогда $\exists z_1, z_2 (f(z_1) = f(z_2))$, не умаляя общности можно считать, что $f(z_1) = f(z_2) = 0$, иначе вычтем константу. Выберем такой $\varepsilon$, что $z_1 \in \overline{B_\varepsilon(z_1)} \subset G$ и $z_2 \in \overline{B_\varepsilon(z_2)} \subset G$, и $B_\varepsilon(z_1) \inter B_\varepsilon(z_2) = \varnothing$. По утверждению (\ref{uccs}), $f$ аналитична, значит, можно написать
    \begin{gather*}
        f(z) = 0 + (z - z_1)^{k_1} g_1(z), \quad g_1(z_1) \neq 0
    \end{gather*}
    Тогда при $z \in \partial B_\varepsilon(z_1)$
    \begin{equation*}
        |f(z)| = \varepsilon^{k_1} \cdot |g_1(z)| = \varepsilon^{k_1} \cdot |g_1(z_1) + o(1)| = \varepsilon^{k_1} \cdot |g_1(z_1)| + o(\varepsilon^{k_1})
    \end{equation*}
    Последнее выражение при достаточно малом $\varepsilon$ строго больше $\frac{\varepsilon^{k_1}}{2} g_1(z_1)$. $\partial B_\varepsilon(z_1)$ --- компакт, значит, на нём $f_n \rightrightarrows f$. Тогда для достаточно большого $N$ для любого $n > N$ верно
    \begin{equation*}
        |f| > \frac{\varepsilon^{k_1}}{2} g_1(z_1) > \frac{\varepsilon^{k_1}}{4} g_1(z_1) > |f_n - f|
    \end{equation*}

    Теперь можно применить теорему Руше (\ref{rouche}) для пары функций $(f, f_n - f)$ и контура $\partial B_\varepsilon(z_1)$. Получаем, что у функций $f$ и $f + (f_n - f) = f_n$ одинаковое число корней внутри $\partial B_\varepsilon(z_1)$. У $f$ такой корень хотя бы 1 (собственно, $z_1$). Таким образом, для достаточно большого $N$ при любом $n > N$ $f_n$ имеет корень в $B_\varepsilon(z_1)$.

    По абсолютно аналогичному рассуждению, $f_n$ имеет корень также и в $B_\varepsilon(z_2)$. $B_\varepsilon(z_1) \inter B_\varepsilon(z_2) = \varnothing$, следовательно, эти корни не совпадают, значит, у $f_n$ хотя бы 2 корня. Но если $f_n$ --- однолистная, то это невозможно по определению.
\end{proof}

\begin{definition}[Локальная однолистность]
    Функция $f$ называется \textit{локально однолистной} в $G$, если вокруг всякой точки $z_0 \in G$ найдётся окрестность, в которой $f$ однолистна.
\end{definition}

\begin{theorem}[О локальной однолистности]
    $f$ локально однолистна $\iff \forall z_0 \in G (f'(z_0) \neq 0)$. 
\end{theorem}
\begin{proof}
    $\Ra$: Предположим противное, то есть нашлась точка $z_0$ такая, что $f'(z_0) = 0$. Тогда
    \begin{equation*}
        f(z) = f(z_0) + (z - z_0)^k g(z),\, k \geq 2,\, g(z_0) \neq 0
    \end{equation*}

    Раз уж $g(z_0) \neq 0$, можно выбрать $\varepsilon$ такой, что в $B_\varepsilon(z_0)$ $|g(z)| \geq a$. В $B_\varepsilon(z_0)$ $g(z) \neq 0$, значит, можно определить $h(z)$ такое, что $e^{h(z)} = g(z)$. Обозначим $\tilde{g}(z) = e^{h(z) / k}$. Тогда
    \begin{equation*}
        f(z) = f(z_0) + \left( (z - z_0) \tilde{g}(z) \right)^k
    \end{equation*}

    Будем решать уравнение $f(z) = f(z_0) + w$ для фиксированного $w$ малого модуля. Это эквивалентно уравнению $(z - z_0)\tilde{g}(z) = s \in \sqrt[k]{w}$. Заметим, что если уравнение такого вида можно решить вне зависимости от выбора конкретного $s$, то, поскольку таких $s$ ровно $k \geq 2$, функция $f(z)$ окажется не однолистной.

    Обозначим $\tilde{g}(z) = \tilde{g}(z_0) + \tilde{\tilde{g}}(z)$, где $\tilde{\tilde{g}}(z) = o(1)$. К функциям $(z - z_0)\tilde{g}(z_0) - s$, $\tilde{\tilde{g}}(z)(z - z_0)$ можно применить теорему Руше (\ref{rouche}): \todo{почему}? Функция $(z - z_0)\tilde{g}(z_0) - s$ линейна, то есть у неё есть ровно 1 корень - тогда у функции $(z - z_0)(\tilde{g}(z_0) + \tilde{\tilde{g}}(z)) - s$ тоже есть 1 корень, что и требовалось.

    $\La$: не умаляя общности, $f(z_0) = 0$ (иначе можно вычесть константу). Знаем, что $f'(z_0) \neq 0$. Тогда
    \begin{equation}\tag{$\star$}
        f(z) = f'(z_0)(z - z_0) + g(z), \quad g(z) = o(z - z_0)
    \end{equation}

    Пусть $f(z_1) = f(z_2)$, $|z - z_1|, |z - z_2| < \varepsilon$. Подставим $z_1, z_2$ в формулу выше и вычтем:
    \begin{equation*}
        f'(z_0)(z_1 - z_2) + g(z_1) - g(z_2) = f(z_1) - f(z_2) = 0
    \end{equation*}
    \begin{equation*}
        |g(z_1) - g(z_2)| = \left| \int_{z_1}^{z_2} g'(s)ds \right| \leq |z_1 - z_2| \underset{|z_0 - s| < \varepsilon}{\max} |g'(s)|
    \end{equation*}

    Покажем, что этот $\max$ настолько быстро стремится к 0 при $\varepsilon \to 0$, что равенство $(\star)$ не может выполняться.
    \begin{equation*}
        \underset{|z_0 - s| < \varepsilon}{\max} |g'(s)| \leq \underset{|z_0 - s| < \varepsilon}{\max} \left| \frac{1}{2 \pi i} \int_{S_\delta(z)} \frac{g(w)}{(w - z)^2} dw \right| \leq \frac{2 \pi \delta}{2 \pi} \cdot \frac{1}{\delta^2} \cdot \underset{|z_0 - s| < \varepsilon}{\max} |g(s)| = \frac{1}{\delta} \underset{|z_0 - s| < \varepsilon}{\max} |g(s)|
    \end{equation*}
    Это верно для всякого $\delta$. Пусть $z_1, z_2 \in B_{\frac{\varepsilon}{2}}(z_0)$. Расмотрим $\delta = \frac{\varepsilon}{2}$:
    \begin{equation*}
        \frac{2}{\varepsilon} \underset{|z_0 - s| < \varepsilon}{\max} |g(s)| \underset{\varepsilon \to 0}{\longrightarrow} 0
    \end{equation*}
    так как $g(z) = o(z - z_0)$. Подставляя итоговую оценку в $(\star)$, перенося $g(z_1) - g(z_2)$ в правую часть, и беря модуль, получаем
    \begin{equation*}
        |f'(z_0) (z_1 - z_2)| \leq |z_1 - z_2| o(1)
    \end{equation*}
    Левая часть убывает линейно и не равна 0, а вторая убывает быстрее при $\varepsilon \to 0$.
\end{proof}

\begin{definition}[Компактность в $A(G)$]
    Семейство $\{f_n\}_{n \in I} \subset A(G)$ называется \textit{компактным}, если в любой последовательсти функций из него содержится сходящаяся в $A(G)$ последовательность.
\end{definition}
\begin{remark}
    Вообще говоря, это свойство правильно называть "<относительная секвенциальнциальная компактность">, но в этом курсе мы будет просто говорить "<компактность">.
\end{remark}

\begin{lemma}
    Если $\{f_n\}_{n \in I}$ равномерно ограничена, то и $\{f_n'\}_{n \in I}$ равномерно ограничена.
\end{lemma}
\begin{proof}
    \begin{equation*}
        f_n'(z) = \frac{1}{2 \pi i} \int_{S_\varepsilon(z)} \frac{f_n(w)}{(w - z)^2} dw
    \end{equation*}
\end{proof}

\begin{theorem}[Монтеля]
\label{montel}
    $\{f_n\}_{n \in I}$ - компактна $\iff$ $\{f_n\}_{n \in I}$ равномерно ограничена, то есть
    \begin{equation*}
        \forall K \subset G, \text{K - компакт } \exists C = C(K) \forall z \in K \forall n \in I \big( |f_n(z)| \leq C(K) \big)
    \end{equation*}
\end{theorem}
\begin{proof}
    $\La$: Зафиксируем компакт $K$. Рассмотрим в $K$ конечную $\varepsilon$-сеть. 
    
    \todo{вспомнить что тут было вообще}
    \begin{equation*}
        \forall \varepsilon > 0 \exists \delta = \delta(\varepsilon, K) \forall z_1, z_2 \in K \forall f_n \in \{f_n\}_{n \in I} \big( |z_1 - z_2| < \delta \Ra |f_n(z_1) - f_n(z_2)| < \varepsilon \big)
    \end{equation*}

    $\Ra$: Предположим противное, то есть $\exists K \left( \underset{z \in K}{\sup} |f_n(z)| \underset{n \to +\infty}{\longrightarrow} +\infty \right)$. Можно считать $|f_l(z_l)| > l$, тогда сходящейся подпоследовательности не существует.
\end{proof}

\begin{theorem}[Арцела-Асколи]
    $K$ --- метрический компакт. Тогда $\{f_i\} \subset C(K)$ компактна $\iff$ $\{f_i\}$ равномерно ограничена и равностепенно непрерывна.
\end{theorem}

\begin{definition}[Конформное отображение]
    \todo{.}
\end{definition}

\begin{theorem}[Римана]
    $G$ --- односвязная область в $\hat{\C}$, $|\partial G| \geq 2$. Тогда $G$ переводится в единичный открытый диск $D$ конформным отображением.
\end{theorem}
\begin{proof}
    Рассмотрим класс функций $\{\varphi\} = I$ такой что
    \begin{enumerate}
        \item $\varphi \in A(G)$
        \item $\varphi$ --- однолистная
        \item $\varphi(G) \subset D$
    \end{enumerate}

    Будем максимизировать функционал $J(\varphi) = |\varphi'(a)|$, где $a$ --- некоторая фиксированная точка в $G$.

    Сначала, добавим в $I$ константные функции, по модулю не превосходящие 1. Теперь $I$ --- замкнутый класс по теореме об однолистности (\ref{one_leaf}).

    Во-вторых, заметим, что $J$ --- непрерывный функционал, поскольку
    \begin{equation*}
        \varphi'(a) = \frac{1}{2 \pi i} \int_{S_\varepsilon(a)} \frac{\varphi(w)}{(w - a)^2} dw
    \end{equation*}
    \todo{и что?}

    Значит, по теореме Арцела-Асколи (\todo{как её применить?}) $I$ --- компакт, значит, $A = \underset{\varphi \in I}{\sup} |\varphi'(a)|$ достигается. Покажем, что $A \neq 0$.

    Если есть $z_0 \notin \operatorname*{Cl}(G)$, то функция $\frac{1}{z - z_0}$ ограничена на $G$, а значит, для достаточно малого $\varepsilon > 0$ функция $\frac{\varepsilon}{z - z_0} \in I$, поскольку модуль её значений можно сделать $< 1$, а однолисточно очевидна.
    
    Однако, может оказаться, что $\operatorname*{Cl}(G) = \hat{\C}$ (например, $G = \hat{\C} \setminus [0, 1]$). Рассмотрим две различные точки $\alpha, \beta \in \partial G$. Поскольку $G$ односвязна, можно однозначно определить функцию $\sqrt{\frac{z - \alpha}{\beta}}$:
    \begin{equation*}
        g(z) = \frac{z - \alpha}{z - \beta} \Ra h(z) = \sqrt{\frac{z - \alpha}{z - \beta}} = e^{\frac{1}{2} \log(g(z))}
    \end{equation*}
    При этом таких функций две, $h_1(z), h_2(z)$, поскольку у $\sqrt{}$ две ветви. Каждая функция $h_i$ однолистна:
    \begin{equation*}
        h_i(z_1) = h_i(z_2) \Ra \frac{z_1 - \alpha}{z_1 - \beta} = h_i^2(z_1) = h_i^2(z_2) = \frac{z_2 - \alpha}{z_2 - \beta} \Ra (z_1 - \alpha)(z_2 - \beta) = (z_2 - \alpha)(z_1 - \beta) \Ra z_1 = z_2
    \end{equation*}
    Естественным образом, $h_1(z) = -h_2(z)$. Оказывается, что их образы не пересекаются:
    \begin{equation*}
        h_1(z_1) = h_2(z_2) \Ra h_1(z_1) = -h_1(z_2) \Ra h_1^2(z_1) = h_1^2(z_2) \Ra z_1 = z_2
    \end{equation*}
    (последнее следствие проводится аналогично предыдущей формуле). 

    $h_i$ есть конформное отображение на $h_i(G)$, значит, \todo{можно сделать} аналогично первому случаю.

    Итого, мы выяснили, что $A \neq 0$. Теперь покажем, что функция, на которой достигается этот супремум, должна действовать во весь $D$. Обозначим эту функцию $\varphi_0$.

    Заметим, что $\varphi_0(a) = 0$. Предположим противное, $\varphi_0(a) \neq 0$. Тогда
    \begin{gather*}
        \xi(z) = \frac{\varphi_0(z) - \varphi_0(a)}{1 - \overline{\varphi_0(a)} \varphi_0(z)} = \mu_{\varphi_0(a)} \circ \varphi_0 \\
        \xi' = \frac{\varphi_0'(a)}{1 - \overline{\varphi_0(a)} \varphi_0(a)} = \frac{\varphi_0'(a)}{1 - |\varphi_0(a)|^2} \\
        1 - |\varphi_0(a)|^2 \in \R, \leq 1
    \end{gather*}
    Но тогда $\varphi_0$ --- не максимум $J$, потому что $J(\xi) > J(\varphi)$, противоречие.

    Пусть $\varphi_0(G) \neq D$, $b \in D \setminus \varphi_0(G)$.
    \todo{Дописать доказательство}.
\end{proof}

\end{document}